\chapter*{Vorwort: Was Regeln nicht leisten}
\addcontentsline{toc}{chapter}{Vorwort}
\markboth{Vorwort}{Vorwort}

\textbf{Dreiundzwanzig Bände haben Regeln aufgestellt.} \kursiv{Dieser prüft sie an zwanzig Fällen --- und in sechs davon lautet die Antwort, dass man nichts ändert. Das ist der eigentliche Gegenstand: nicht die Regel, sondern die Entscheidung, ob sie hier gilt.}

\begin{merksatz}
\boxtitel{Warum ein eigener Band}
\begin{center}
\renewcommand{\arraystretch}{1.5}
\begin{tabular}{@{}P{2.8cm}P{2.4cm}P{2.8cm}@{}}
\toprule
& \textbf{Eine Regel} & \textbf{Ein Fall}\\
\midrule
Sagt & \textbf{was meistens gilt} & \textbf{was hier gilt}\\
Enthält & Verallge\-meinerung & \textcolor{gruen}{\textbf{Umstände}}\\
Anwendbar & \textcolor{akzent}{ohne Urteil} & \textcolor{gruen}{mit Urteil}\\
Falsch, wenn & \textcolor{akzent}{\textbf{die Umstände abweichen}} & \\
\bottomrule
\end{tabular}
\end{center}

\begin{center}
\begin{minipage}{0.88\textwidth}
\kursiv{Eine Regel ohne Umstände ist ein Dogma. Die Bände 17 bis 23 haben zu jeder Regel den Anwendungsbereich genannt --- aber wer eine Regel liest, merkt sich die Regel und vergisst den Bereich. Ein Fall enthält beides untrennbar.}
\end{minipage}
\end{center}

\textbf{Der Aufbau jedes Falls.}

\begin{center}
\renewcommand{\arraystretch}{1.5}
\begin{tabular}{@{}P{2.6cm}P{6.2cm}@{}}
\toprule
\textbf{Die Vorlage} & Der Code, wie er vorliegt, mit seinem Zusammenhang\\
\textbf{Die Diagnose} & \textbf{Was daran auffällt --- und was nicht}\\
\textbf{Das Urteil} & \textbf{Überarbeiten, umbauen oder lassen}\\
\textbf{Die Begründung} & Warum, mit Zahlen wo möglich\\
\textbf{Der Fehlschluss} & \textbf{Was man daraus \kursiv{nicht} lernen sollte}\\
\bottomrule
\end{tabular}
\end{center}

\begin{center}
\kursiv{Der fünfte Abschnitt ist der ungewöhnliche und der wichtigste. Zu jedem richtigen Urteil gehört eine naheliegende Verallgemeinerung, die falsch ist --- und wer nur das Urteil mitnimmt, nimmt sie mit.}
\end{center}
\end{merksatz}

\begin{warnung}
\boxtitelr{Die drei Ausgänge}
\begin{center}
{\small
\renewcommand{\arraystretch}{1.5}
\begin{tabular}{@{}P{2.4cm}P{1.6cm}P{4.4cm}@{}}
\toprule
\textbf{Urteil} & \textbf{Fälle} & \textbf{Bedeutet}\\
\midrule
\textbf{Überarbeitet} & $\mathbf{9}$ & \textbf{derselbe Entwurf, besser ausgeführt}\\
\textbf{Umgebaut} & $\mathbf{5}$ & \textbf{der Entwurf war falsch}\\
\textbf{So gelassen} & $\mathbf{6}$ & \textcolor{gruen}{\textbf{die Vorlage war richtig}}\\
\bottomrule
\end{tabular}}
\end{center}

\medskip
\begin{center}
\begin{minipage}{0.88\textwidth}
\kursiv{Sechs von zwanzig bleiben, wie sie sind --- und in vier davon widerspricht das einer Regel aus einem früheren Band. Das ist kein Widerruf, sondern die Anwendung derselben Regel mit ihrem Anwendungsbereich.}
\end{minipage}
\end{center}

\begin{center}
\kursiv{Wer einen Band über Urteilskraft schreibt, in dem jeder Fall geändert wird, hat einen Katalog geschrieben. Die Fälle 15 bis 20 sind deshalb der Kern dieses Bandes, nicht sein Anhang.}
\end{center}
\end{warnung}

\begin{merksatz}
\boxtitel{Woher die Fälle kommen}
\begin{center}
{\small
\renewcommand{\arraystretch}{1.45}
\begin{tabular}{@{}P{2.4cm}P{2.2cm}P{2.4cm}P{1.8cm}@{}}
\toprule
\textbf{Teil} & \textbf{Fälle} & \textbf{Gegenstand} & \textbf{Bände}\\
\midrule
I & $1$--$4$ & \textbf{Die Daten} & $20$, $23$\\
II & $5$--$9$ & \textbf{Der Ablauf} & $21$, $22$\\
III & $10$--$14$ & Die Struktur & $17$, $20$\\
IV & $15$--$20$ & \textbf{Wann man es lässt} & alle\\
\bottomrule
\end{tabular}}
\end{center}

\begin{center}
\begin{minipage}{0.88\textwidth}
\kursiv{Alle Fälle stammen aus demselben Vorhaben, das die Bände 20 bis 23 begleitet hat --- Messerfassung, Übertragung, Speicherung, Auswertung. Sie sind zugespitzt, aber nicht erfunden: Jeder beruht auf einer Bauform, die in klinischen Systemen vorkommt.}
\end{minipage}
\end{center}

\textbf{Die Sprachen.}

\begin{center}
\renewcommand{\arraystretch}{1.5}
\begin{tabular}{@{}P{2.4cm}P{2.0cm}P{3.4cm}@{}}
\toprule
\textbf{Sprache} & \textbf{Fälle} & \\
\midrule
Python & $7$ & \textbf{Band 20}\\
C & $5$ & Band 21\\
Go & $5$ & Band 22\\
Sprachfrei & $3$ & \textbf{Nähte, Band 23}\\
\bottomrule
\end{tabular}
\end{center}

\begin{center}
\kursiv{Die letzte Zeile ist wichtig: Drei Fälle enthalten überhaupt keinen Code, sondern eine Vereinbarung. Sie gehören zu den folgenreichsten, weil eine falsche Vereinbarung länger lebt als jedes Programm.}
\end{center}
\end{merksatz}

\section*{Wie man diesen Band liest}

\textbf{Erstens: die Vorlage beurteilen, bevor man die Diagnose liest.}

\begin{center}
\begin{minipage}{0.88\textwidth}
\kursiv{Jeder Fall beginnt mit dem Code und seinem Zusammenhang. Wer an dieser Stelle innehält und sein eigenes Urteil bildet, hat den Band gewinnbringend gelesen --- auch dann, wenn sein Urteil vom hier vertretenen abweicht.}
\end{minipage}
\end{center}

\textbf{Zweitens: das Urteil ist nicht die Antwort.}

Es gibt Fälle, in denen ein anderes Urteil ebenso vertretbar ist. Wo das so ist, steht es dabei. Was nicht verhandelbar ist, ist die \kursiv{Begründung} --- ein Urteil ohne Begründung ist Geschmack.

\textbf{Drittens: einzeln lesen.}

\begin{center}
\kursiv{Zwanzig Fälle hintereinander erzeugen den Eindruck, überall sei etwas faul. Das ist der Fehler jedes Fallbuchs. Ein Fall je Sitzung, mit eigenem Urteil vorweg, ist die vorgesehene Verwendung.}
\end{center}

\begin{warnung}
\boxtitelr{Drei Warnungen}
\textbf{Erstens: Kein Fall ist eine Regel.}

\begin{center}
\begin{minipage}{0.88\textwidth}
\kursiv{Aus \enquote{hier war die Abstraktion verfrüht} folgt nicht \enquote{Abstraktion ist verfrüht}. Der Abschnitt \enquote{Der Fehlschluss} steht in jedem Fall genau deshalb dort.}
\end{minipage}
\end{center}

\textbf{Zweitens: Die Zahlen sind Größenordnungen.}

Verstehenszeiten, Laufzeiten, Aufwände --- alle gemessen oder geschätzt, alle mit Streuung. Die Verhältnisse tragen, die absoluten Werte nicht.

\textbf{Drittens: Der Bestand ist nicht schlecht.}

\begin{center}
\kursiv{Jede Vorlage in diesem Band lief im Betrieb und erfüllte ihren Zweck. Es geht nicht darum, Fehler vorzuführen, sondern darum, Entscheidungen sichtbar zu machen, die beim Schreiben unsichtbar waren.}
\end{center}
\end{warnung}

\section*{Zum Aufbau}

\begin{itemize}
\item \textbf{Teil I} (Fälle 1--4): Die Daten. Einheit, Argumente, \textbf{verfrühte Abstraktion}, das fehlende Feld.
\item \textbf{Teil II} (Fälle 5--9): Der Ablauf. Schleife, \textbf{entfernte Prüfung}, Puffer, Goroutine, Kanal.
\item \textbf{Teil III} (Fälle 10--14): Die Struktur. Hierarchie, Zwischenspeicher, \textbf{Zeit}, Wiederholung, Fehlermeldung.
\item \textbf{Teil IV} (Fälle 15--20): \textbf{Wann man es lässt.} Sechs Fälle ohne Änderung.
\end{itemize}

\begin{geschichte}
\boxtitel{Warum \enquote{Die Kunst}}
Der Titel ist eine Anleihe und eine Behauptung.

\begin{center}
\renewcommand{\arraystretch}{1.5}
\begin{tabular}{@{}P{2.6cm}P{6.2cm}@{}}
\toprule
\textbf{Die Anleihe} & Knuth nannte sein Werk \enquote{The Art of Computer Programming} --- nicht Wissenschaft, nicht Technik\\
\textbf{Die Begründung} & \textbf{Weil das Urteil nicht aus dem Wissen folgt}\\
\textbf{Die Behauptung} & Man kann es dennoch lehren --- an Fällen, nicht an Regeln\\
\bottomrule
\end{tabular}
\end{center}

\begin{center}
\begin{minipage}{0.86\textwidth}
\kursiv{Zwischen \enquote{ich kenne die Regel} und \enquote{ich weiß, ob sie hier gilt} liegt eine Lücke, die kein Lehrbuch schließt. Die Medizin kennt sie und hat eine Antwort gefunden, die dreitausend Jahre alt ist: die Falldarstellung.}
\end{minipage}
\end{center}

\textbf{Was die Falldarstellung leistet.}

\begin{center}
\renewcommand{\arraystretch}{1.5}
\begin{tabular}{@{}P{2.8cm}P{2.4cm}P{2.6cm}@{}}
\toprule
\textbf{Eigenschaft} & \textbf{Regel} & \textbf{Fall}\\
\midrule
Merkbar & \textcolor{gruen}{leicht} & \textcolor{gruen}{\textbf{leichter}}\\
Übertragbar & \textcolor{gruen}{scheinbar} & \textcolor{mittel}{mit Vorsicht}\\
Enthält Umstände & \textcolor{akzent}{\textbf{nein}} & \textcolor{gruen}{\textbf{ja}}\\
Widerspruch möglich & \textcolor{akzent}{nein} & \textcolor{gruen}{\textbf{ja}}\\
\bottomrule
\end{tabular}
\end{center}

\begin{center}
\begin{minipage}{0.86\textwidth}
\kursiv{Die letzte Zeile ist der Grund für die Form: Über eine Regel kann man nur zustimmen oder ablehnen. Über einen Fall kann man streiten --- und im Streit über einen konkreten Fall lernt man mehr als im Streit über eine Regel.}
\end{minipage}
\end{center}

\begin{center}
\kursiv{Deshalb enthält dieser Band keine neuen Regeln. Er enthält zwanzig Gelegenheiten, die vorhandenen anzuwenden --- und sechs Gelegenheiten, sie nicht anzuwenden.}
\end{center}
\end{geschichte}

\bigskip
\begin{flushright}
\itshape\color{grau}
Regeln sagen, was meistens gilt.\\
Die Kunst ist zu wissen,\\
ob dies einer der meisten Fälle ist.\\[10pt]
\normalfont\color{grau}\upshape
--- der Leitgedanke dieses Bandes.\\
\itshape Wer keine Regeln kennt, kann nicht urteilen.\\
Wer nur Regeln kennt, urteilt nicht.
\end{flushright}
